\documentclass{beamer}
\usetheme{Madrid}
\usecolortheme{default}

\usepackage[utf8]{inputenc}
\usepackage[T1]{fontenc}
\usepackage[french]{babel}
\usepackage{booktabs}
\usepackage{graphicx}

\definecolor{lacunared}{HTML}{E53935}
\definecolor{lacunablue}{HTML}{1E88E5}
\definecolor{lacunadark}{HTML}{14161C}

\setbeamercolor{palette primary}{bg=lacunadark, fg=white}
\setbeamercolor{palette secondary}{bg=lacunablue, fg=white}
\setbeamercolor{palette tertiary}{bg=lacunared, fg=white}
\setbeamercolor{titlelike}{parent=palette primary}
\setbeamercolor{structure}{fg=lacunablue}

\title{Présentation Technique du Projet Lacuna}
\subtitle{Architecture MVC, IA, Réseau TCP \& Validation IHM}
\author{Groupe 14}
\date{\today}

\begin{document}

\begin{frame}
    \titlepage
\end{frame}

\begin{frame}{Membres du Projet et Encadrement}
    \begin{columns}[t]
        \begin{column}{0.5\textwidth}
            \begin{block}{Membres du Groupe}
                \begin{itemize}
                    \item MADJOUR Amir
                    \item FRITAH Wassim
                    \item ZOUITEN Ouassim
                    \item BOUGUERRA Lakhdar
                    \item DJERBOA Mohamed
                    \item SAHRAOUI Mohammed
                \end{itemize}
            \end{block}
        \end{column}
        \begin{column}{0.48\textwidth}
            \begin{block}{Encadrement}
                \begin{itemize}
                    \item \textbf{Encadreur :} \\ Benjamin Wack
                \end{itemize}
            \end{block}
            \vspace{0.2cm}
            \begin{block}{Contexte}
                Projet de programmation \\
                L3 Informatique -- UGA \\
                Module PROG6
            \end{block}
        \end{column}
    \end{columns}
\end{frame}

\begin{frame}{Sommaire / Agenda}
    \tableofcontents
\end{frame}

\section{Introduction et Concept du Jeu}

\begin{frame}{Le Jeu Lacuna : Concept et Objectif}
    \begin{itemize}
        \item \textbf{Concept :} Adaptation numérique d'un jeu de plateau abstrait et géométrique.
        \item \textbf{Plateau :} Un espace circulaire continu en 2D.
        \item \textbf{Éléments :} \textbf{49 fleurs} réparties au hasard, appartenant à \textbf{7 couleurs distinctes} (7 de chaque couleur).
        \item \textbf{But du jeu :} Capturer des fleurs pour obtenir la majorité (au moins 4 fleurs sur 7) dans au moins \textbf{4 couleurs différentes}.
    \end{itemize}
    \begin{block}{Condition de Victoire}
        Le premier joueur à cumuler au moins 4 majorités de couleurs de fleurs est déclaré vainqueur.
    \end{block}
\end{frame}

\begin{frame}{Déroulement du Jeu : Les Deux Phases}
    \begin{enumerate}
        \item \textbf{La phase de placement et capture directe :}
        \begin{itemize}
            \item Chaque joueur dispose de 6 pions.
            \item À tour de rôle, placement d'un pion sur le segment de droite reliant deux fleurs de même couleur.
            \item Capture instantanée et retrait des deux fleurs cibles du plateau.
            \item \textit{Contraintes :} La ligne ne doit traverser aucun obstacle (pion, fleur) et le pion ne doit pas être trop proche d'une fleur.
        \end{itemize}
        \item \textbf{La phase de proximité finale (Hold-up) :}
        \begin{itemize}
            \item Après le placement des 12 pions, les fleurs restantes sont attribuées au joueur possédant le pion le plus proche (calcul de distance minimale).
        \end{itemize}
    \end{enumerate}
\end{frame}

\begin{frame}{Problématiques et Enjeux du Projet}
    \begin{itemize}
        \item \textbf{Enjeu Géométrique :} Gérer les calculs de segments de droite, de hitboxes et de distances minimales en virgule flottante sur un espace 2D continu.
        \item \textbf{Enjeu d'Interaction :} Traduire des gestes physiques du jeu de société vers une IHM informatique (visée, aimantation).
        \item \textbf{Enjeu d'IA :} Créer une intelligence artificielle compétitive capable de simuler à la fois la capture directe et d'anticiper la phase de Voronoï.
        \item \textbf{Enjeu Réseau :} Implémenter un protocole d'échange de données léger et résistant aux déconnexions pour le mode multijoueur.
    \end{itemize}
\end{frame}

\section{Architecture Logicielle}

\begin{frame}{Architecture Globale : Modèle-Vue-Contrôleur (MVC)}
    \begin{block}{Le patron de conception MVC au cœur du projet}
        \begin{description}
            \item[Modèle (\texttt{GameModel}) :] Gère les structures de données pures. Il maintient la liste des fleurs, des pions placés, des joueurs, l'historique des coups pour l'annulation/rétablissement (Undo/Redo) et l'état de la partie.
            \item[Vue (\texttt{MainFrame}, \texttt{BoardPanel}, \texttt{PlayerPanel}) :] Affiche graphiquement les données du modèle à l'aide de Java Swing et réagit aux évènements graphiques (souris, boutons).
            \item[Contrôleur (\texttt{GameController}) :] Interprète les clics de l'utilisateur, vérifie la validité des coups selon les règles et orchestre le passage des tours (Joueur, IA ou Joueur distant).
        \end{description}
    \end{block}
\end{frame}

\begin{frame}{Diagramme Conceptuel des Classes}
    \begin{center}
        \begin{tabular}{ccc}
            \toprule
            \textbf{Modèle} & \textbf{Vue} & \textbf{Contrôleur} \\
            \midrule
            \texttt{GameModel} & \texttt{MainFrame} & \texttt{GameController} \\
            \texttt{Player} & \texttt{BoardPanel} & \\
            \texttt{Pawn} & \texttt{PlayerPanel} & \\
            \texttt{Flower} & \texttt{BoardRenderer} & \\
            \texttt{FlowerColor} & \texttt{TutorialOverlay} & \\
            \bottomrule
        \end{tabular}
    \end{center}
    \begin{itemize}
        \item Utilisation du motif de conception \textbf{Observer/Observable} : le modèle notifie ses auditeurs (\texttt{ModelListener}) lors de toute mise à jour de son état pour déclencher un rafraîchissement visuel (\texttt{repaint()}).
    \end{itemize}
\end{frame}

\section{Conception et Interactions IHM}

\begin{frame}{Conception de l'IHM : L'Expérience Visuelle}
    \begin{itemize}
        \item \textbf{Thème Sombre :} Fond principal gris foncé (\texttt{\#14161C}) pour un meilleur confort visuel et un contraste saisissant avec les pions-fleurs lumineux.
        \item \textbf{Identité Visuelle Joueurs :} Attribution de couleurs claires et distinctives (Rouge électrique \texttt{\#E53935} pour le Joueur 1, Bleu intense \texttt{\#1E88E5} pour le Joueur 2).
        \item \textbf{Symétrie Épurée :} Plateau central sous forme d'ellipse, encadré par des compteurs latéraux affichant en temps réel les majorités acquises.
    \end{itemize}
\end{frame}

\begin{frame}{Interaction Utilisateur \& Assistance au Placement (Feedforward)}
    \begin{itemize}
        \item \textbf{Pion Fantôme :} Le pointeur système par défaut est masqué au survol du tapis et remplacé par une prévisualisation semi-transparente du pion du joueur.
        \item \textbf{Lignes d'alignement :} Dès que le pointeur passe entre deux fleurs de même couleur admissibles, le jeu dessine la ligne reliant les fleurs.
        \item \textbf{Résolution des conflits (La Molette) :} Si plusieurs lignes valides se croisent sous le pointeur, un popup contextuel invite à tourner la \textbf{molette de la souris} pour alterner simplement entre les paires candidates.
    \end{itemize}
\end{frame}

\begin{frame}{Retours Visuels \& Animations (Feedback)}
    \begin{itemize}
        \item \textbf{Cylinder Intro :} Animation au démarrage mimant l'éjection et l'atterrissage des fleurs sur le tapis de feutre, agrémentée de cercles d'ondes de choc.
        \item \textbf{Vol des fleurs :} Lors de la capture, les fleurs tournoient physiquement depuis le plateau pour aller s'empiler dans le tableau de scores latéral.
        \item \textbf{Boutons d'action rapides (Undo/Redo) :} Permettent de revenir en arrière ou de rétablir un coup à l'aide d'alertes textuelles temporisées ("Toasts").
    \end{itemize}
\end{frame}

\section{Intelligence Artificielle}

\begin{frame}{Intelligence Artificielle : Modélisation de l'État}
    \begin{itemize}
        \item L'IA fonctionne grâce à une représentation allégée et optimisée de l'état du jeu via la classe \texttt{AIState}.
        \item Un générateur de mouvements (\texttt{DefaultAIMoveGenerator}) explore à chaque étape de la phase de placement l'ensemble des paires de fleurs admissibles sur le plateau.
        \item Pour chaque paire, il calcule le point de placement optimal du pion (généralement le centre géométrique, ajusté si des obstructions sont détectées).
    \end{itemize}
\end{frame}

\begin{frame}{L'Algorithme Minimax et Élagage Alpha-Beta}
    \begin{block}{Algorithme de décision}
        Recherche arborescente basée sur l'algorithme \textbf{Minimax} permettant de simuler ses propres coups (maximisation) et les réponses de l'adversaire (minimisation).
    \end{block}
    \begin{itemize}
        \item \textbf{Élagage Alpha-Beta :} Coupe les branches de l'arbre dont l'évaluation est moins bonne que la meilleure valeur déjà garantie, réduisant drastiquement le nombre de configurations explorées.
        \item \textbf{Table de Transposition :} Mémorise les états déjà évalués dans une table de hachage pour éviter de recalculer des branches redondantes.
    \end{itemize}
\end{frame}

\begin{frame}{Fonction d'Évaluation de l'IA}
    La classe \texttt{DefaultAIEvaluator} estime la force d'une position selon :
    \begin{enumerate}
        \item \textbf{Le score de capture directe :} Nombre de fleurs déjà engrangées.
        \item \textbf{La proximité géométrique (Voronoï) :} Évaluation dynamique des distances entre les pions et les fleurs restantes. L'IA simule l'attribution de fin de partie pour savoir si un placement sécurise une zone géographique.
        \item \textbf{La valeur des majorités :} Pondération spéciale pour favoriser l'obtention de la $4^e$ fleur d'une couleur (seuil critique de majorité).
    \end{enumerate}
\end{frame}

\begin{frame}{Fonction d'Évaluation : Optimisation des Paramètres}
    \begin{itemize}
        \item \textbf{\texttt{TunableAIEvaluator} :} Permet de modifier à la volée les coefficients d'importance (poids de la capture immédiate contre poids de la proximité finale).
        \item \textbf{\texttt{MultiParamSweepRunner} :} Outil d'entraînement automatisé faisant s'affronter des IA configurées avec différentes combinaisons de poids pour extraire la paramétrisation statistiquement la plus forte.
    \end{itemize}
\end{frame}

\section{Implémentation Réseau}

\begin{frame}{Implémentation Réseau : Architecture Globale du Serveur}
    \begin{block}{Serveur TCP Multithreadé (\texttt{Server.java})}
        \begin{itemize}
            \item \textbf{Sockets standard :} Utilisation de \texttt{ServerSocket} sur le port \texttt{38400} par défaut pour des liaisons fiables.
            \item \textbf{Gestion des clients :} Un thread d'écoute accepte les connexions entrantes et lance immédiatement un thread dédié de type \texttt{Connexion\_handler} pour chaque joueur.
        \end{itemize}
    \end{block}
    \begin{block}{Interface d'Administration (\texttt{Server\_command\_line.java})}
        \begin{itemize}
            \item Lancement d'un thread console indépendant de l'écoute réseau.
            \item Commandes disponibles : \texttt{HELP} et \texttt{LIST} pour afficher en temps réel l'état des salons (ID, joueurs, visibilité, temps d'inactivité).
        \end{itemize}
    \end{block}
\end{frame}

\begin{frame}{Gestion des Sessions et Salons de Jeu}
    \begin{block}{Le Gestionnaire de Sessions (\texttt{SessionManager.java})}
        \begin{itemize}
            \item Utilisation d'un dictionnaire concurrentiel thread-safe (\texttt{ConcurrentHashMap}) pour stocker et manipuler les sessions actives.
            \item Génération automatique de clés uniques de \textbf{5 caractères alphanumériques} via un générateur aléatoire sécurisé (\texttt{SecureRandom}).
        \end{itemize}
    \end{block}
    \begin{block}{Types de Visibilité}
        \begin{itemize}
            \item \textbf{Sessions publiques :} Listées dans le lobby et accessibles à tous les joueurs en attente.
            \item \textbf{Sessions privées (\texttt{prv}) :} Accessibles uniquement en saisissant manuellement le code unique du salon.
        \end{itemize}
    \end{block}
\end{frame}

\begin{frame}{Cycle de Vie des Salons et Robustesse}
    \begin{block}{Maintien et Nettoyage de la Mémoire}
        \begin{itemize}
            \item \textbf{Mécanisme anti-fuite :} Un horodatage de la dernière action (\texttt{lastTouchedTime}) est mis à jour à chaque interaction client.
            \item \textbf{Expiration automatique :} Les salons inactifs depuis plus de 30 minutes sont automatiquement marqués comme obsolètes (\texttt{isOutdated}) et purgés lors des requêtes d'accès.
        \end{itemize}
    \end{block}
    \begin{block}{Fermeture Propre}
        \begin{itemize}
            \item Suppression dynamique de la session de la table de hachage lorsque le salon se vide.
            \item Gestion de la fermeture ordonnée des flux d'E/S et de la socket cliente.
        \end{itemize}
    \end{block}
\end{frame}

\begin{frame}{Protocole d'Appairage (Handshake et Lobby)}
    Le protocole repose sur des commandes textuelles simples échangées en mode texte :
    \begin{description}
        \item[\texttt{PING} :] Test de disponibilité de la connexion, le serveur répond \texttt{pong}.
        \item[\texttt{LIST} :] Demande la liste des salons publics libres. Le serveur renvoie le nombre de salons suivi de paires \texttt{<nom\_hote> <id\_salon>}.
        \item[\texttt{CREAT <nom> [prv]} :] Crée une nouvelle session. Le serveur répond \texttt{OK <id\_session>} en cas de succès.
        \item[\texttt{JOIN <nom> <id>} :] Permet à un visiteur de rejoindre une session spécifique. Le serveur notifie l'hôte via \texttt{JOINED <nom>} et valide avec \texttt{OK}.
    \end{description}
\end{frame}

\begin{frame}{Le Relais Universel de Données (Forwarding)}
    Une fois la session d'appairage validée, le serveur agit comme un pont de transport transparent (\texttt{forwardMessages}) :
    \begin{block}{Le Conteneur \texttt{DATA <n>}}
        \begin{itemize}
            \item Les clients envoient une enveloppe contenant la commande \texttt{DATA} suivie du nombre de lignes à transmettre \texttt{n}.
            \item Le serveur lit les \texttt{n} lignes suivantes et les redirige immédiatement vers l'adversaire (peer) sans chercher à décoder leur structure.
        \end{itemize}
    \end{block}
    \begin{block}{Avantages de ce Choix de Conception}
        \begin{itemize}
            \item \textbf{Découplage absolu} entre la logique métier du jeu (positions des fleurs, coups joués) et l'architecture réseau.
            \item \textbf{Évolutivité facilitée} : modification ou enrichissement du protocole du jeu sans impact sur le code du serveur.
        \end{itemize}
    \end{block}
\end{frame}

\begin{frame}{Gestion des Erreurs et Robustesse en Partie}
    \begin{block}{Codes de Statut du Serveur}
        \begin{description}
            \item[\texttt{BAD 400} :] Requête malformée ou syntaxe incorrecte.
            \item[\texttt{BAD 404} :] Commande inconnue.
            \item[\texttt{BAD 405} :] Salon de jeu déjà complet.
            \item[\texttt{BAD 500} :] Salon introuvable (ID erroné ou expiré).
        \end{description}
    \end{block}
    \begin{block}{Déconnexions Impromptues}
        \begin{itemize}
            \item Si une connexion socket est coupée pendant la partie, le serveur envoie automatiquement la commande \texttt{LEFT} à l'autre joueur.
            \item Permet au client de réagir proprement (affichage d'une boîte de dialogue de déconnexion et retour sécurisé au menu principal).
        \end{itemize}
    \end{block}
\end{frame}

\section{Validation de l'IHM et Tests Utilisateurs}

\begin{frame}{Validation IHM : Protocole d'Évaluation}
    \begin{itemize}
        \item \textbf{Objectif :} Évaluer l'ergonomie, la clarté (feedforward/feedback) et l'accessibilité de l'interface.
        \item \textbf{Panel :} 12 participants répartis en 3 groupes cibles :
        \begin{itemize}
            \item \textbf{Groupe 1 : Gamers} (Mohammed, Ilyes, Radhi, Nadir)
            \item \textbf{Groupe 2 : Non-Gamers} (Marouan, Djaber, Khaled, Sami)
            \item \textbf{Groupe 3 : Seniors} (La mère d'Amir, Le père d'Amir, L'oncle de Wassim, La mère de Wassim)
        \end{itemize}
        \item \textbf{Méthodologie :} Session de jeu complète suivie d'un entretien semi-directif à base de questions/réponses.
    \end{itemize}
\end{frame}

\begin{frame}{Retours d'Évaluation par Panel}
    \begin{itemize}
        \item \textbf{Gamers :} apprécient la réactivité du pion fantôme, l'affichage de Voronoï, mais demandent de pouvoir désactiver les intros lentes et d'avoir des raccourcis clavier.
        \item \textbf{Non-Gamers :} plébiscitent l'overlay interactif du tutoriel qui permet d'apprendre sans lire de longs règlements. La molette de sélection est jugée très naturelle.
        \item \textbf{Seniors :} soulignent le grand confort du thème sombre pour les yeux et l'aide à la visée (aimantation), mais regrettent la petite taille des caractères des scores latéraux.
    \end{itemize}
\end{frame}

\begin{frame}{Améliorations Ergonomiques Apportées}
    Suite aux retours des 12 participants, trois ajustements majeurs ont été intégrés :
    \begin{enumerate}
        \item \textbf{Aimantation adaptative :} Augmentation de la sensibilité d'accroche du pion fantôme sur les lignes pour pallier le manque de précision ou les tremblements physiques.
        \item \textbf{Raccourcis clavier :} Implémentation de \texttt{Ctrl+Z} (Annuler) et \texttt{Ctrl+Y} (Rétablir) demandés par les gamers.
        \item \textbf{Temporisation de transition :} Adoucissement de l'enchaînement avec la phase de Voronoï afin de ne pas surprendre l'utilisateur.
    \end{enumerate}
\end{frame}

\section{Méthodologie de Travail et Organisation du Groupe}

\begin{frame}{Méthodologie de Travail \& Organisation du Groupe}
    \begin{block}{Méthode Agile Scrum}
        \begin{itemize}
            \item \textbf{Scrum Master :} MADJOUR Amir
            \item Organisation en sprints réguliers
        \end{itemize}
    \end{block}
    \begin{block}{Outils Collaboratifs}
        \begin{itemize}
            \item \textbf{Gestion de tâches :} ClickUp
            \item \textbf{Suivi de version (code) :} Git
            \item \textbf{Communication :} WhatsApp \& Discord (travail d'équipe et revues)
        \end{itemize}
    \end{block}
    \begin{block}{Coordination et Revues Post-Audit}
        \begin{itemize}
            \item \textbf{Assemblage obligatoire} en visioconférence après chaque audit ou rendez-vous avec le tuteur (Benjamin Wack) pour traiter les retours et aligner l'équipe.
        \end{itemize}
    \end{block}
\end{frame}

\section{Conclusion \& Perspectives}

\begin{frame}{Conclusion \& Perspectives}
    \begin{itemize}
        \item \textbf{Bilan :} Validation réussie. L'IHM de Lacuna parvient à concilier la rigueur de sa physique géométrique continue avec une ergonomie intuitive et accessible.
        \item \textbf{Perspectives :}
        \begin{itemize}
            \item Mode accessibilité "Grandes polices" pour les seniors.
            \item Schéma de symboles géométriques complémentaires pour les joueurs daltoniens.
            \item Portage de l'interface Swing vers des architectures mobiles tactiles.
        \end{itemize}
    \end{itemize}
\end{frame}

\begin{frame}
    \begin{center}
        \Huge Merci pour votre attention !\\
        \vspace{1cm}
        \large Questions ?
    \end{center}
\end{frame}

\end{document}
